% 共同版式与课程约束；在生成的 report_template/ 中可按教师要求调整。
\usepackage{amsmath,amssymb,graphicx,booktabs,array,tabularx,float}
\usepackage{geometry,xcolor,caption,siunitx,fancyhdr,tikz}
\usepackage[super,sort&compress]{gbt7714}
\setcitestyle{super,sort&compress}
\usepackage{hyperref}
\hypersetup{hidelinks}
\geometry{left=2cm,right=2cm,top=2cm,bottom=2cm,headheight=15pt,headsep=6mm}
\setlength{\parindent}{2em}
\setlength{\parskip}{0pt}
\setlength{\columnsep}{7mm}
\setlength{\textfloatsep}{10pt plus 2pt minus 2pt}
\setlength{\intextsep}{10pt plus 2pt minus 2pt}
\setlength{\floatsep}{8pt plus 2pt minus 2pt}
\setcounter{topnumber}{3}
\renewcommand{\topfraction}{0.9}
\renewcommand{\textfraction}{0.1}
\renewcommand{\floatpagefraction}{0.75}
\renewcommand{\baselinestretch}{1.12}
\sisetup{detect-all,per-mode=symbol,separate-uncertainty=true}
\definecolor{ReportAccent}{HTML}{222222}
\newcommand{\ReportTitleFont}{\heiti\fontsize{18pt}{23pt}\selectfont}
\newcommand{\ReportTitleAlign}{\centering}
\newcommand{\ReportRule}{\par\vspace{3pt}\hrule\vspace{6pt}}
\newcommand{\ReportHeader}{北京师范大学物理实验报告}
\ifcase\ReportTemplateID\relax
\or % 01：保留 thuemp 双栏，使用统一的两行基本信息。
  \renewcommand{\baselinestretch}{1.15}
  \renewcommand{\ReportHeader}{北京师范大学普通物理实验}
\or % 02：居中衬线标题、紧凑摘要与细线，参考 APS 两栏成稿。
  \renewcommand{\ReportTitleFont}{\songti\bfseries\fontsize{17pt}{22pt}\selectfont}
  \setlength{\columnsep}{6mm}
  \renewcommand{\baselinestretch}{1.10}
\or % 03：左对齐黑体大标题、较疏栏距与蓝色标题层级。
  \renewcommand{\ReportTitleAlign}{\raggedright}
  \renewcommand{\ReportTitleFont}{\heiti\fontsize{21pt}{26pt}\selectfont}
  \setlength{\columnsep}{8mm}
  \definecolor{ReportAccent}{HTML}{163A5F}
  \renewcommand{\ReportRule}{\par\vspace{4pt}{\color{ReportAccent}\hrule height 0.8pt}\vspace{7pt}}
\or % 04：紧凑题头、窄栏距与小节间距，保留课程所需完整讨论。
  \renewcommand{\ReportTitleFont}{\heiti\fontsize{16pt}{20pt}\selectfont}
  \setlength{\columnsep}{5.5mm}
  \renewcommand{\baselinestretch}{1.08}
\or % 05：左对齐题头、开放留白、灰色双细线。
  \renewcommand{\ReportTitleAlign}{\raggedright}
  \renewcommand{\ReportTitleFont}{\songti\bfseries\fontsize{19pt}{25pt}\selectfont}
  \setlength{\columnsep}{7.5mm}
  \renewcommand{\baselinestretch}{1.13}
  \renewcommand{\ReportRule}{\par\vspace{4pt}{\color{black!55}\hrule\vspace{2pt}\hrule}\vspace{7pt}}
\fi
\ctexset{
  section={name={,、},number=\chinese{section},format=\large\heiti\color{ReportAccent},beforeskip=12pt,afterskip=6pt},
  subsection={format=\normalsize\heiti,beforeskip=8pt,afterskip=4pt}
}
\ifnum\ReportTemplateID=4
  \ctexset{section={beforeskip=9pt,afterskip=4pt}}
\fi
\captionsetup{font=small,labelfont=bf,labelsep=quad,skip=5pt}
\captionsetup[figure]{name=图}
\captionsetup[table]{name=表}
\pagestyle{fancy}
\fancyhf{}
\fancyhead[L]{\footnotesize\ReportHeader}
\fancyhead[R]{\footnotesize 实验报告}
\fancyfoot[C]{\footnotesize\thepage}
\renewcommand{\headrulewidth}{0.3pt}
\fancypagestyle{plain}{\fancyhf{}\fancyfoot[C]{\footnotesize\thepage}\renewcommand{\headrulewidth}{0pt}}
\newcommand{\ReportBlank}[1]{\makebox[#1]{\hrulefill}}
\newcommand{\ReportTODO}[1]{\textbf{待补充：}#1}
\newcommand{\ReportEnglishSummary}{%
  \par\bigskip
  \begin{center}\bfseries Abstract\end{center}
  \ReportAbstractEN\par
  \noindent\textbf{Keywords:}\ \ReportKeywordsEN\par}
\newif\ifReportDraft
\ReportDrafttrue
\newif\ifReportAppendix
\ReportAppendixfalse
% 校徽左边界 = 正文左边界向左 2cm；等比例宽7cm，仅首页。
\AddToHook{shipout/foreground}{%
  \ifnum\value{page}=1\relax
    \IfFileExists{assets/bnu_logo.png}{%
      \begin{tikzpicture}[remember picture,overlay]
        \node[anchor=north west,inner sep=0pt,xshift=0cm,yshift=-0.6cm]
          at (current page.north west) {\includegraphics[width=7cm]{assets/bnu_logo.png}};
      \end{tikzpicture}%
    }{\PackageWarning{report-template}{Missing assets/bnu_logo.png; course logo must be supplied}}%
  \fi
}
\makeatletter
\newcommand{\MakeReportFrontMatter}{%
  \twocolumn[{\begin{@twocolumnfalse}
    \vspace*{8mm}
    {\ReportTitleAlign\ReportTitleFont\ReportTitle\par}
    \vspace{5pt}
    \begin{center}
      作者（姓名）：\ReportAuthor\quad 学号：\ReportStudentID\\[4pt]
      指导老师：\ReportTeacher\quad 实验时间：\ReportExperimentDate
    \end{center}
    \ifReportDraft
    \noindent\fcolorbox{black!25}{black!3}{\parbox{\dimexpr\textwidth-2\fboxsep-2\fboxrule\relax}{%
      \small 模板草稿：所有“待补充”均须以讲义、原始记录或可核实资料填写；占位图表与文献不能作为实测结果提交。
    }}\par\vspace{6pt}
    \fi
    \noindent\textbf{摘要}\quad\ReportAbstractCN\par
    \noindent\textbf{关键词：}\ReportKeywordsCN\par\vspace{6pt}
    {\ReportTitleAlign\bfseries\ReportTitleEN\par}
    {\ReportTitleAlign\small\ReportAuthorEN\par}\vspace{3pt}
    \ReportRule
  \end{@twocolumnfalse}}]
  \thispagestyle{plain}
}
\makeatother
